% 整数（综述）
% license CCBYSA3
% type Wiki

本文根据 CC-BY-SA 协议转载翻译自维基百科\href{https://en.wikipedia.org/wiki/Integer}{相关文章}

整数是指 零（0）、正自然数（1, 2, 3, …），或 正自然数的相反数（−1, −2, −3, …）。\(^\text{[1]}\)正自然数的相反数（或加法逆元）称为负整数。\(^\text{[2]}\)所有整数构成的集合通常记作 **粗体 Z** 或 黑板粗体：$\mathbb{Z}$.\(^\text{[3][4]}\)  

自然数集合  $\mathbb{N}$是整数集合$\mathbb{Z}$的子集，而 $\mathbb{Z}$ 又是有理数集合$\mathbb{Q}$的子集，而 $\mathbb{Q}$ 本身是实数集合$\mathbb{R}$的子集。\(^\text{[a]}\)与自然数集合类似，整数集合$\mathbb{Z}$是可数无限集。整数可以被视为一种实数，其表示中不含分数部分。例如：$21, 4, 0, -2048$ 都是整数；而 $9.75, 5+\tfrac{1}{2}, \tfrac{5}{4}, \sqrt{2}$ 都不是整数。\(^\text{[7]}\)

整数构成了包含自然数的最小的群与最小的环。在代数数论中，整数有时被称为有理整数，以区别于更一般的代数整数。实际上，（有理）整数就是那些既是代数整数、又是有理数的数。
